\chapter{部署、运行与维护}
\label{ch:operations}

\section{环境与构建}
项目兼容 ROS 2 Foxy/Jazzy：Edge 设备保持 Foxy，开发/仿真环境可使用 Jazzy。首次使用仿真时按 \pathref{scripts/setup_workspace_python.sh} 创建工作区 Python 环境；AI 依赖根据 GPU/CUDA 条件选择性安装。两个 workspace 应按 overlay 顺序构建：先 \pkg{workspace_auv}，再 \pkg{workspace_sim}。

\begin{lstlisting}[language=bash,caption={构建和环境检查}]
cd workspace_auv
colcon build --symlink-install
source install/setup.bash
python3 -m py_compile $(find src -name '*.py')

cd ../workspace_sim
colcon build
source install/setup.bash
\end{lstlisting}

\section{真机启动前检查}
\begin{enumerate}
  \item 确认电池、电源、急停、推进器机械状态和舱体密封。
  \item 确认 Jetson/Edge 已连接 ROS 2 网络，MCU 通道或外部 Agent 可用。
  \item 确认相机设备路径、分辨率、时间戳和标定 profile。
  \item 确认 ZIT6 状态、INS 导航就绪、heartbeat 和电池状态。
  \item 先运行无运动的 readiness 检查，再启用 motion/task。
  \item 下水前确认 Mission 文件、目标类别、速度上限和超时策略。
\end{enumerate}

\section{运行观测}
核心观测包括 Topic 是否存在、消息频率、节点列表、任务状态、估计状态、ZIT6 状态、控制周期、相机 MJPEG/annotated 流和日志目录。\pkg{uv_log} 可以同时记录 rosbag、视频片段和节点日志；停止时应允许 recorder 完成索引和尾部恢复。

\section{部署}
仓库提供基于 rsync 的部署脚本，默认目标为 AUV 电脑上的项目目录。首次部署推荐使用 \code{--dry-run}，确认后再同步；选择性部署修改文件时使用 checksum 或 Git changed 文件列表。\code{--delete} 会删除远端多余文件，只能在确认本地和远端目录边界后使用。

\section{关闭流程}
先停止任务和运动，再关闭预览/记录，确认控制器发送安全状态，最后关闭 launch 和 ROS 2。若发生异常退出，优先保存日志和 rosbag，使用 \pkg{uv_log recover} 恢复完整片段后再清理结果目录。

\section{长期维护}
代码、接口、参数和文档应在同一提交中变更；每个正式实验绑定 Git commit；每次 protocol 或坐标系变更都要更新契约测试；新算法必须说明输入 Topic、输出 Topic、频率、延迟、失败行为和真值隔离情况。发布前执行 \code{git diff --check}、Python 编译检查、接口测试和对应模式的启动 smoke test。

\begin{warningbox}
真机 profile 与仿真 profile 不可混用。任何会改变布防、最大速度、推力或 watchdog 的配置变更，都必须经过静态检查、台架测试和有急停条件的水池测试。
\end{warningbox}

